% N5: bounded-query (poly-query) hardness theorem -- DRAFT.
% LaTeX fragment, not a standalone document.  Companion to
%   results/N5_asymptotics.py     (sympy, output results/N5_asymptotics.txt, 15/15)
%   results/N5_delta_at_etahat.py (LP oracle re-run, output
%                                  results/N5_delta_at_etahat.txt, 40/0/2)
% and it quotes, without re-deriving, the oracle facts of
%   RESEARCH_STATE.md R9 / R11(a)  and  results/T2_summary.md (Conclusion 1 and 3),
%   results/N4_hardness_construction.md (Section 3 and Section 6),
%   RESEARCH_STATE.md R1 (the greedy guarantee L_K) and R7/T5 (the U_K family).
%
% Conventions (paper + CLAUDE.md):
%   alpha in (0,1], larger is better;  d_e(S) = f(S cup {e}) - f(S);
%   eta_u = max d/dtilde, eta_o = max dtilde/d, eta = eta_u eta_o;
%   SINGLE-ELEMENT error model throughout this file (see vacuity check V7);
%   L_K(eta) = 1 - (1 - 1/(eta K))^K,  U_K(eta) = 1 - (1 - 1/(eta(K-1)+1))^K.
%
% Status tags per CLAUDE.md.  Nothing in this file is claimed as "proved":
% every deduction step is [HAND-PROOF-UNREVIEWED], every instance-level fact is
% imported with the tag it carries in RESEARCH_STATE.md, the asymptotic and
% monotonicity statements are [VERIFIED-SYMBOLIC] via results/N5_asymptotics.py,
% and the error-inflation constant is [VERIFIED-LP] at 40 points via
% results/N5_delta_at_etahat.py.
%
% NEW IN N5 (not a restatement of T2): Lemma N5.4 CORRECTS the closed form of
% results/T2_summary.md Conclusion 1 / RESEARCH_STATE.md R11(a); the LP oracle
% refutes the published form at 8 of 40 points.  See Lemma N5.4 and D1(b).
%
% SCOPE WORD (GLOSSARY.md, "any algorithm"): everything below lives in range
% (ii) poly-query = hidden O + concentration.  Range (i) (unbounded queries) is
% R2 and gives the DIFFERENT constant 1/eta; range (iii) (poly-time) is not
% touched anywhere in this paper.


% ===========================================================================
% 0.  The instance family
% ===========================================================================
% Status: the family itself is the R9 candidate, restated verbatim from
% RESEARCH_STATE.md R9 / results/T2_hardness_grid.py (function fixedF_min_delta).
% [VERIFIED-LP, finite instances] for monotonicity, submodularity, OPT = F(O),
% the ratio identity and the error band -- see Lemma N5.4 and the pointers there.
%
% Ground set N with |N| = n.  Hidden O subseteq N with |O| = K, B = N \ O.
% For S subseteq N put  x = |S \ O|,  y = |S cap O|.
% Design parameter theta >= 1 and a_theta = 1 - 1/(theta K).  Below, the
% instantiation used by the theorem is written hateta and hata = a_hateta;
% Lemmas N5.4 and N5.5 are stated for a general theta because the theorem has to
% SOLVE for the right one.
% Balancedness threshold tau with 1 <= tau < K.
%
%   f  = F_O(S)  = hateta^{-1/2} * ( 1 - hata^x (1 - y/K) )
%   ftilde = G_O(S) = 1 - hata^{x+y}                      if y <= tau
%                   = 1 - hata^{x+tau} (K - y)/(K - tau)  if y >= tau
%   Ghat(s) = 1 - hata^s        (the balanced-region profile; G_O(S) = Ghat(|S|)
%                                whenever y <= tau, since then x + y = |S|)
%
% The two branches agree at y = tau.  The error split is
% hateta_u = hateta_o = sqrt(hateta); the prefactor hateta^{-1/2} on F only fixes
% that split and cancels from every ratio below.
% NOTE (used later): f(O) = F_O(O) = hateta^{-1/2}, and F_O(S) <= f(O) for all
% |S| <= K, i.e. OPT = f(O) -- [VERIFIED-LP] by full enumeration, T2 step 3.


% ===========================================================================
% 1.  Concentration lemma
% ===========================================================================
% Status: [HAND-PROOF-UNREVIEWED].  Elementary union bound + hypergeometric
% ratio; no oracle was run for it.  No script dependency.
\begin{lemma}[Balancedness of a fixed small query]\label{lem:n5-conc}
Let $O$ be a uniformly random $K$-subset of a ground set $N$ with $|N|=n\ge K$,
let $S\subseteq N$ be fixed with $|S|\le K$, and let $0\le\tau<K$ be an integer.
Then
\[
  \Pr\bigl[\,|S\cap O|>\tau\,\bigr]\;\le\;\binom{K}{\tau+1}\Bigl(\frac{K}{n}\Bigr)^{\tau+1}
  \;\le\;\frac{K^{2(\tau+1)}}{(\tau+1)!\;n^{\tau+1}} .
\]
\end{lemma}
% Proof draft.  Write m = tau + 1.  The event {|S cap O| > tau} is the union,
% over the C(|S|, m) subsets R subseteq S with |R| = m, of the events {R subseteq O}.
% For a fixed such R,
%       Pr[R subseteq O] = C(n-m, K-m)/C(n, K) = prod_{i=0}^{m-1} (K-i)/(n-i),
% and (K-i)/(n-i) <= K/n for every 0 <= i < m whenever K <= n (cross-multiplying,
% (K-i)n <= K(n-i)  <=>  iK <= in).  Hence Pr[R subseteq O] <= (K/n)^m.
% A union bound over the C(|S|, m) <= C(K, m) <= K^m/m! choices of R gives the claim.
% [此步需人工检查：C(|S|,m) <= C(K,m) 用了 |S| <= K；m <= K 保证二项式系数非零。]

% ---------------------------------------------------------------------------
% Status: [HAND-PROOF-UNREVIEWED] (arithmetic on Lemma N5.1 only).
\begin{lemma}[Union bound over a transcript, and an explicit $n$]\label{lem:n5-union}
Fix an integer $c\ge 0$, set $Q=n^{c}$ and $\tau=c+1$, and let $K\ge\tau+1$.  Let
$S_1,\dots,S_Q$ and $T_0$ be fixed sets with $|S_i|\le K$ and $|T_0|\le K$, and
let $O$ be a uniformly random $K$-subset of $N$.  Then
\[
  \Pr\Bigl[\,\exists i:\ |S_i\cap O|>\tau\ \ \text{ or }\ \ T_0\cap O\neq\emptyset\,\Bigr]
  \;\le\;\frac{K^{2c+4}}{(c+2)!\;n^{2}}\;+\;\frac{K^{2}}{n}\;<\;\tfrac12
\]
whenever $n\ge 4K^{\,c+2}$ (and $K\ge 2$).
\end{lemma}
% Proof draft.  First term: Lemma N5.1 with tau + 1 = c + 2 gives, per query,
% C(K, c+2)(K/n)^{c+2} <= K^{2c+4}/((c+2)! n^{c+2}); multiplying by Q = n^c gives
% K^{2c+4}/((c+2)! n^2).  Second term: Pr[T_0 cap O != empty] <= E|T_0 cap O|
% = |T_0| K/n <= K^2/n (Markov, or a union bound over the elements of T_0; the
% hypergeometric mean E|T_0 cap O| = |T_0| K / n is exact).
% Numerically, n >= 4 K^{c+2} gives n^2 >= 16 K^{2c+4}, so the first term is at
% most 1/(16 (c+2)!) <= 1/32, and the second is at most K^2/(4K^{c+2}) =
% 1/(4K^c) <= 1/4; the sum is at most 9/32 < 1/2.
% [此步需人工检查：(c+2)! >= 2 用了 c >= 0；K^c >= 1 用了 K >= 1。]
% REMARK (vacuity check V11 below): tau = c is already enough for the union
% bound (it gives K^{2c+2}/((c+1)! n) instead of the 1/n^2 above), and tau enters
% delta linearly, so tau = c would give a slightly STRONGER theorem.  tau = c+1
% is kept because TASKS2.md fixes it and because tau = 0 is degenerate at c = 0.


% ===========================================================================
% 2.  Valuation lemma
% ===========================================================================
% Status: [HAND-PROOF-UNREVIEWED] for the simulation/induction argument.
% It CONSUMES four oracle facts about the family of Section 0, all from
% RESEARCH_STATE.md R11(a) = results/T2_summary.md Conclusion 1 and step 3
% ([VERIFIED-LP] on the finite test set listed in honest declaration D1,
%  [CONJECTURE] for general (K, tau, eta)):
%
%  (F1) balanced answers depend only on |S|:  G_O(S) = Ghat(|S|) whenever
%       |S cap O| <= tau.  (This is a definitional property of the candidate G
%       of R9; in the LP it is the balanced-region equality constraint.)
%  (F2) F_O is monotone submodular, and the pair (F_O, G_O) has single-element
%       error exactly hateta(1 + delta(hateta)) with 1 + delta(theta) =
%       max{a_theta^tau K/(K-tau), a_theta^{1-tau}}^2 -- Lemma N5.4 below.
%  (F3) OPT = max_{|S| <= K} F_O(S) = F_O(O)  (full enumeration, T2 step 3;
%       re-checked on the grid at the theorem's hateta values by
%       results/N5_delta_at_etahat.py, column C).
%  (F4) F_O(T)/F_O(O) = 1 - hata^{K} = L_K(hateta) for any K-set T with T cap O
%       = emptyset (T2 step 3: the columns ratio_B and LK of results/T2_table.csv
%       agree row by row; re-checked at the hateta values, column D).
\begin{lemma}[Valuation on the balanced event]\label{lem:n5-val}
Let $\mathcal A$ be a deterministic algorithm that makes at most $Q$ queries to
$\tilde f$, each on a set of size at most $K$, and returns a set $T$ with
$|T|\le K$.  Run $\mathcal A$ against the \emph{canonical} oracle that answers
every query $S$ with $\hat G(|S|)$; let $S_1,\dots,S_Q$ be the resulting queries
and $T_0$ the resulting output.  Both are fixed sets, independent of $O$.  If
$O$ satisfies
\[
  |S_i\cap O|\le\tau\ \ (1\le i\le Q)\qquad\text{and}\qquad T_0\cap O=\emptyset,
\]
then $\mathcal A$, run on the instance $(F_O,G_O)$, outputs exactly $T_0$ and
\[
  \frac{F_O(T_0)}{F_O(O)}\;=\;1-\hat a^{\,|T_0|}\;\le\;1-\hat a^{K}
  \;=\;1-\Bigl(1-\frac{1}{\hat\eta K}\Bigr)^{K}\;=\;L_K(\hat\eta).
\]
\end{lemma}
% Proof draft.  Induction on i.  Suppose the first i-1 real answers agree with
% the canonical ones; then the algorithm is deterministic and its state after
% i-1 answers is the canonical state, so its i-th query is the canonical S_i.
% By hypothesis |S_i cap O| <= tau, so by (F1) G_O(S_i) = Ghat(|S_i|), which is
% the canonical answer.  Hence all Q answers agree and the output is T_0.
% Since T_0 cap O = empty we have y = 0 and x = |T_0| <= K, so
%   F_O(T_0) = hateta^{-1/2}(1 - hata^{|T_0|}) <= hateta^{-1/2}(1 - hata^{K}),
% using hata in (0,1); and F_O(O) = hateta^{-1/2}, which is OPT by (F3).  The
% last equality is (F4).
% [此步需人工检查：(a) "算法是确定性的，故状态由应答序列唯一决定" 这一步排除了
%  算法用随机串或用 O 的其他信息；(b) 询问集合大小 <= K 是 Lemma N5.1 的前提，
%  这里必须与算法模型一致；(c) hata in (0,1) 需要 hateta K > 1。]
%
% VARIANT used by the randomized corollary (same proof, weaker conclusion):
% if only the query condition |S_i cap O| <= tau holds (T_0 cap O may be
% nonempty), then the output is still T_0 and, since F_O(x,y) is
% hateta^{-1/2}(1 - hata^x + hata^x y/K) <= hateta^{-1/2}(1 - hata^{K}) +
% hateta^{-1/2} y/K for x <= K,
%     F_O(T_0)/F_O(O) <= L_K(hateta) + |T_0 cap O| / K .
% Taking expectations over O and using E|T_0 cap O| = |T_0| K/n <= K^2/n gives
% the additive term K/n of Corollary N5.7.
% [此步需人工检查：hata^x <= 1 与 x <= K 处的单调性方向。]


% ===========================================================================
% 3.  The error-inflation constant delta
% ===========================================================================
% *** THIS SECTION CORRECTS results/T2_summary.md CONCLUSION 1. ***
% Status of the corrected closed form (the max below): [VERIFIED-LP] at 40
%   (K, eta, tau) points, results/N5_delta_at_etahat.py -> N5_delta_at_etahat.txt
%   (40 PASS / 0 FAIL / 2 SKIP), on top of the T2 points; [CONJECTURE] in general.
% Status of the derivation below (which band constraint binds):
%   [HAND-PROOF-UNREVIEWED]; it reproduces the quantity delta_cc computed by
%   results/T2_hardness_grid.py and now matches the LP minimum at 8 further
%   points where the T2 formula does NOT.
% Status of the first-order dominance condition theta >= 2 - 1/tau:
%   [VERIFIED-SYMBOLIC], results/N5_asymptotics.py check (iii).
\begin{lemma}[Error of the family]\label{lem:n5-delta}
For the family of Section~0 with design parameter $\theta>1$ and
$a_\theta=1-1/(\theta K)$, let $s\ge 1$ be the smallest number such that the
inflated single-element band
$\Delta F/(\sqrt\theta\, s)\le \Delta G\le \sqrt\theta\, s\,\Delta F$ holds on
every adjacent pair of sets.  Then
\[
  1+\delta(\theta)\;=\;s^{2}\;=\;
  \max\Bigl\{\ a_\theta^{\tau}\frac{K}{K-\tau}\ ,\ a_\theta^{\,1-\tau}\ \Bigr\}^{2},
\]
and the pair $(F_O,G_O)$ then has single-element error exactly
$\theta\,(1+\delta(\theta))$.
\end{lemma}
% Proof draft of the lower bound (the four edge computations).  Both endpoints
% balanced, eta_u = eta_o = sqrt(theta), 1 - a = 1/(theta K), write a = a_theta:
%   x-edge (x,y) -> (x+1,y), y <= tau:
%       Delta F = theta^{-1/2} a^x (1-a)(1 - y/K),  Delta G = a^{x+y}(1-a)
%       Delta G/(eta_o Delta F) = a^{y} K/(K-y)          <- first candidate
%       Delta F/(eta_u Delta G) = (K-y)/(theta K a^{y}) <= 1   (uses theta >= 1)
%   y-edge (x,y) -> (x,y+1), y+1 <= tau:
%       Delta F = theta^{-1/2} a^x / K,  Delta G = a^{x+y}(1-a)
%       Delta G/(eta_o Delta F) = a^{y}/theta <= 1             (uses theta >= 1)
%       Delta F/(eta_u Delta G) = a^{-y}                <- second candidate
% The map y |-> a^{y} K/(K-y) is increasing on 0 <= y <= tau (the ratio of
% consecutive terms is a (K-y)/(K-y-1) >= 1 iff a >= 1 - 1/(K-y), which holds
% since a = 1 - 1/(theta K) and theta >= 1), so its max over the balanced strip
% is at y = tau; a^{-y} is maximized at y = tau - 1, giving a^{1-tau}.  These
% are "constant-constant" constraints: both endpoints lie in the balanced region
% where G is pinned to Ghat, so no choice of the unbalanced part of G can relax
% them.  Sufficiency (an extension of G to the unbalanced region exists at this
% s, and the explicit G of Section 0 is one) is the LP part.
% [此步需人工检查：四条边比值的代数；"y = tau 的 x-edge 存在" 需要 tau <= K-1
%  且 x+1 <= n-K；两条 "<= 1" 的推导都用了 theta >= 1，见假设 (H)。]
%
% *** CORRECTION TO T2 (this is a new N5 finding, not a restatement). ***
% results/T2_summary.md Conclusion 1 reports 1 + delta = (a^tau K/(K-tau))^2,
% i.e. the FIRST term only.  To first order in 1/K the two candidates are
% tau(1 - 1/theta)/K and (tau-1)/(theta K), so the first dominates iff
% theta >= 2 - 1/tau [VERIFIED-SYMBOLIC, results/N5_asymptotics.py (iii)].
% Every T2 test point satisfies that (T2 used tau <= 2 and theta in {3/2,2,3},
% and 2 - 1/2 = 3/2), which is why the second term was never seen.  The theorem
% below evaluates the family at a design parameter strictly SMALLER than eta,
% which can drop below 2 - 1/tau -- and there the T2 formula is simply wrong.
% [VERIFIED-LP] results/N5_delta_at_etahat.py: 40 points solved, the LP minimum
% delta equals the max-form at all 40, and DIFFERS from the T2 formula at 8 of
% them.  Largest discrepancy (K,eta,tau) = (4,3,2): LP and max-form 0.710334,
% T2 formula 0.367408 (a factor 1.93 in 1 + delta).  Others: (6,1.5,2) 0.404077
% vs 0.141302; (6,2,2) 0.328460 vs 0.274928; (8,1.5,2) 0.247292 vs 0.142723;
% (12,1.5,2) 0.140573 vs 0.106920; (16,1.5,2) 0.098537 vs 0.082317;
% (24,1.5,2) 0.061794 vs 0.055593; (32,1.5,2) 0.045045 vs 0.041807.
% ACTION FOR THE PAPER: T2 Conclusion 1 and RESEARCH_STATE.md R11(a) should be
% amended to the max-form.  (Not edited here: N5 may only write results/N5_*.)

% ---------------------------------------------------------------------------
% Status: [VERIFIED-SYMBOLIC], results/N5_asymptotics.py checks (iv), (vi).
% This lemma REPLACES the naive "delta is increasing in the design parameter"
% argument, which is FALSE for the corrected max-form of Lemma N5.4: the first
% branch increases in theta but the second, a_theta^{1-tau}, DECREASES in it.
\begin{lemma}[The admissibility map is invertible]\label{lem:n5-mono}
Let $K\ge 2\tau$, $\tau\ge1$, and let
$\Phi(\theta):=\theta\,\bigl(1+\delta(\theta)\bigr)$ be the single-element error
of the family with design parameter $\theta$ (Lemma~\ref{lem:n5-delta}).  Then
$\Phi$ is continuous and strictly increasing on $\theta\ge1$, so for every
$\eta\ge\Phi(1)$ there is a unique $\theta^{\ast}=\Phi^{-1}(\eta)\le\eta$; the
family with design parameter $\theta^{\ast}$ has single-element error exactly
$\eta$.
\end{lemma}
% Proof draft.  Phi(theta) = max{ theta c1(theta)^2 , theta c2(theta)^2 } with
% c1 = a^tau K/(K-tau) and c2 = a^{1-tau}, a = a_theta.  A max of increasing
% functions is increasing, so it suffices that both branches increase:
%   theta c1^2 = (K/(K-tau))^2 * theta a^{2 tau}, and
%       d/dtheta [theta a^{2tau}] = a^{2tau} + 2 tau a^{2tau-1}/(theta K) > 0
%       for a > 0                                [VERIFIED-SYMBOLIC, (vi.a)];
%   theta c2^2 = theta a^{2-2tau}, and with u = theta K,
%       d/dtheta [theta a^{2-2tau}] = (1/K)(u/(u-1))^{2tau-2}[1 - (2tau-2)/(u-1)],
%       which is > 0 iff u > 2 tau - 1        [VERIFIED-SYMBOLIC, (vi.b)];
%       theta >= 1 and K >= 2 tau give u >= 2 tau > 2 tau - 1.
% Phi(theta) >= theta because delta >= 0, hence Phi^{-1}(eta) <= eta.
% [此步需人工检查：delta >= 0 需要 max{c1, c2} >= 1，由 c1 >= 1 或 c2 >= 1 得到
%  （tau = 1 时 c2 = 1）。K >= 2 tau 就是定理里的假设 (H2)。]
%
% WHY A FIXED POINT AND NOT AN EXPLICIT FORMULA (honest declaration D3).
% TASKS2.md prescribes the explicit hateta = eta/(1 + delta(eta)).  That value
% is admissible (family error <= eta) iff delta(hateta) <= delta(eta), which is
% exactly the monotonicity that the corrected delta does NOT have in general.
% Two ways out, both recorded:
%   (a) use theta* = Phi^{-1}(eta): always admissible, error EXACTLY eta, and it
%       gives the strongest bound of the family since theta* >= any admissible
%       design parameter and L_K is decreasing (check (v));
%   (b) keep the explicit hateta and CHECK delta(hateta) <= delta(eta).  This was
%       checked at all 40 points of results/N5_delta_at_etahat.py (column E) and
%       held everywhere, so at those points the TASKS2 form is valid -- but it is
%       [VERIFIED-LP finite] there, not proved.
% The theorem below is stated with theta* and records hateta as the explicit
% (weaker, conditionally valid) instantiation.


% ===========================================================================
% 4.  The theorem (deterministic) and its randomized corollary
% ===========================================================================
% Status: [HAND-PROOF-UNREVIEWED] for the assembly (Lemmas N5.1--N5.3, N5.5);
%         [VERIFIED-LP, finite instances] for the facts (F1)--(F4) and for the
%         delta formula; [CONJECTURE] for general (K, tau, theta);
%         [VERIFIED-SYMBOLIC] for the K -> infinity statement in Corollary N5.8.
\begin{theorem}[Bounded-query hardness, deterministic]\label{thm:n5-det}
Fix an integer $c\ge0$, put $\tau=c+1$, and fix a cardinality bound
$K\ge 2\tau$ and an error level $\eta>1$.  For $\theta\ge1$ write
$a_\theta=1-1/(\theta K)$,
\[
  1+\delta(\theta)=\max\Bigl\{a_\theta^{\tau}\tfrac{K}{K-\tau},\,
                              a_\theta^{\,1-\tau}\Bigr\}^{2},
  \qquad \Phi(\theta)=\theta\bigl(1+\delta(\theta)\bigr),
\]
and assume
\[
  \textbf{(H)}\qquad \Phi(1)\le\eta,
  \qquad\text{sufficient: } K\;\ge\;\tau\Bigl(1+\frac{2}{\ln\eta}\Bigr),
\]
so that $\hat\eta:=\Phi^{-1}(\eta)\in[1,\eta]$ is well defined
(Lemma~\ref{lem:n5-mono}).  Let $n\ge 4K^{\,c+2}$ and let $\mathcal A$ be any
deterministic algorithm that makes at most $Q=n^{c}$ queries to $\tilde f$,
each on a set of size at most $K$, and outputs a set $T$ with $|T|\le K$.  Then
there is a set $O\subseteq N$, $|O|=K$, such that the pair
$(f,\tilde f)=(F_O,G_O)$ of Section~0 built with design parameter $\hat\eta$ is
monotone submodular with single-element error exactly $\eta$, has
$O^{\ast}=O$, and
\[
  f(T)\;\le\;\Bigl(1-\bigl(1-\tfrac{1}{\hat\eta K}\bigr)^{K}\Bigr)f(O^{\ast})
  \;=\;L_K(\hat\eta)\,f(O^{\ast}).
\]
\end{theorem}
% Proof draft.  Build the family with design parameter hateta = Phi^{-1}(eta);
% by Lemma N5.5 its single-element error is Phi(hateta) = eta exactly, so every
% member is admissible.  Let S_1, ..., S_Q, T_0 be the canonical transcript of
% Lemma N5.3.  Draw O uniformly at random among K-subsets; by Lemma N5.2 the bad
% event has probability < 1/2 < 1, so SOME O makes all queries balanced and
% misses T_0.  Fix that O.  By Lemma N5.3 the algorithm outputs T_0 and
% F_O(T_0) <= L_K(hateta) F_O(O) = L_K(hateta) OPT.  The probabilistic argument
% is only used to prove existence; the conclusion is about one fixed instance.
% [此步需人工检查：K >= 2 tau 保证 tau < K（否则 delta 无定义）并蕴含 K >= c+2；
%  n >= 4K^{c+2} 同时蕴含 n >= 2K，故 B 至少有 K 个元素、
%  "T_0 cap O = empty 的 K-集" 存在。]
%
% REMARK (recovering the statement prescribed by TASKS2.md).  If additionally
% the first branch of the max attains delta at both eta and at
% hateta_expl := eta/(1 + delta(eta)) -- equivalently a^{2 tau - 1} K >= K - tau
% there, whose first-order form is theta >= 2 - 1/tau [VERIFIED-SYMBOLIC (iii)] --
% then delta(hateta_expl) <= delta(eta), so hateta_expl <= hateta and, since L_K
% is decreasing (check (v)),
%      f(T) <= L_K(hateta) <= L_K(hateta_expl) f(O*)
%           = ( 1 - (1 - 1/(hateta_expl K))^K ) f(O*),
%      hateta_expl = eta/(1+delta),  1 + delta = (a^tau K/(K-tau))^2,  a = 1-1/(eta K),
% which is verbatim the target statement of TASKS2.md N5.  It is a WEAKER
% (larger) bound than the theorem's, and it is conditional on that branch
% condition; the condition was verified at all 40 points of
% results/N5_delta_at_etahat.py (column E).  [VERIFIED-LP finite / CONJECTURE].
%
% WHY HYPOTHESIS (H) IS NOT COSMETIC [HAND-PROOF-UNREVIEWED, found while
% tabulating results/N5_asymptotics.txt and confirmed by
% results/N5_delta_at_etahat.txt].  eta >= 1 always (eta = eta_u eta_o), so a
% design parameter below 1 is meaningless; worse, the two "non-binding" edge
% computations of Lemma N5.4 both use theta >= 1
%   (Delta G/(eta_o Delta F) = a^y/theta <= 1 and
%    Delta F/(eta_u Delta G) = (K-y)/(theta K a^y) <= 1),
% so for theta < 1 further edges can bind and Lemma N5.4 itself fails.
% Concretely: (K, eta, tau) = (4, 3/2, 2) and (4, 2, 2) give
% eta/(1+delta(eta)) = 0.778 and 0.853 -- (H) fails and both are reported
% SKIP(H) in results/N5_delta_at_etahat.txt.  Sufficiency of
% K >= tau(1 + 2/ln eta) for Phi(1) <= eta:
%   ln Phi(1) = ln(1 + delta(1)) <= 2 max{tau ln a + ln(K/(K-tau)), (1-tau) ln a}
%             <= 2 tau/(K - tau) <= ln eta,
% using ln a < 0, -ln(1-t) <= t/(1-t), and (1-tau) ln a <= tau/(K-tau) for
% tau >= 1.
% CONSEQUENCE: for fixed c the theorem is a LARGE-K statement in a second,
% independent sense (V9 gave the first): K must exceed roughly (c+1)(1 + 2/ln eta),
% which blows up as eta -> 1 (see V17).

% ---------------------------------------------------------------------------
% Status: [HAND-PROOF-UNREVIEWED].  Only the AVERAGING (easy) direction of Yao's
% principle is used; no minimax theorem is invoked.  See vacuity check V1.
\begin{corollary}[Bounded-query hardness, randomized]\label{cor:n5-rand}
With the parameters of Theorem~\ref{thm:n5-det}, let $\mathcal R$ be any
randomized algorithm with the same query budget and query-size restriction.
Then
\[
  \min_{O}\ \frac{\mathbb E\bigl[f_O(T)\bigr]}{f_O(O^{\ast})}
  \;\le\;L_K(\hat\eta)\;+\;\varepsilon_n,
  \qquad
  \varepsilon_n=\frac{K}{n}+\frac{K^{2c+4}}{(c+2)!\,n^{2}},
\]
the expectation being over $\mathcal R$'s random string, and the minimum over
the family of Section~0.  In particular $\varepsilon_n\to0$ as $n\to\infty$
with $K,\eta,c$ fixed.
\end{corollary}
% Proof draft.  Let D be the uniform distribution on K-subsets O.  For each
% fixing r of the random string, R_r is deterministic; let T_0^{(r)} be its
% canonical output and E_r the event that all its canonical queries are balanced.
% Bound, pointwise in O, using the VARIANT of Lemma N5.3 and F_O <= F_O(O) (F3):
%    f_O(T)/f_O(O*) <= L_K(hateta) + |T_0^{(r)} cap O|/K + 1{E_r^c}.
% Take E_{O ~ D}: the middle term is at most K/n (E|T_0 cap O| = |T_0|K/n <= K^2/n)
% and the last is at most Q C(K,tau+1)(K/n)^{tau+1} <= K^{2c+4}/((c+2)! n^2)
% (Lemma N5.2, first term only).  Average over r, exchange the two expectations
% (Fubini, finite sums), and use min_O <= E_{O ~ D}.
% [此步需人工检查：1{E_r^c} 处用 f_O(T) <= f_O(O*) 把坏事件的贡献界成 1；
%  以及"最坏实例的期望比 <= 分布平均"这一步的方向。]

% ---------------------------------------------------------------------------
% Status: [VERIFIED-SYMBOLIC], results/N5_asymptotics.py checks (i), (ii), (iii)
% (output: results/N5_asymptotics.txt, 15/15 PASS, exit code 0).
\begin{corollary}[Asymptotics in $K$]\label{cor:n5-asym}
Fix $c$, $\tau=c+1$ and $\eta>1$.  Both branches of $\delta$ vanish as
$K\to\infty$,
\[
  K\,\delta(\eta)\;\longrightarrow\;
  \max\Bigl\{\,2\tau\bigl(1-\tfrac1\eta\bigr),\ \tfrac{2(\tau-1)}{\eta}\,\Bigr\},
  \qquad\text{so}\qquad
  L_K(\hat\eta)\;\longrightarrow\;1-e^{-1/\eta}.
\]
\end{corollary}
% Sympy-verified for symbolic (K, eta, tau) in results/N5_asymptotics.py:
% (i.a) branch-1 delta -> 0, (i.b)/(i.c) K delta -> 2 tau (1 - 1/eta),
% (iii.a)/(iii.b) the two branch first-order constants, (iii.c) branch 1 wins iff
% eta >= 2 - 1/tau, (ii.a) K log(1 - 1/(hateta K)) -> -1/eta, (ii.b) the value
% itself -> 1 - e^{-1/eta}; plus a 30-digit numeric cross-check at K = 10^6.
% Note the CONCLUSION is branch-independent (both branches -> 1, so hateta ->
% eta); only the O(1/K) constant depends on which branch is active.
% Combined with R1 (single-step predictive greedy achieves L_K(eta) ->
% 1 - e^{-1/eta}), the poly-query optimum is pinned to 1 - e^{-1/eta} in the
% limit -- CONDITIONAL on the delta formula, which is [VERIFIED-LP] only on the
% finite set of D1 and [CONJECTURE] in general, and on hypothesis (H), which for
% fixed tau holds for all large K.


% ===========================================================================
% 5.  Self-consistency check (does the theorem contradict R1?)
% ===========================================================================
% Status: [VERIFIED-SYMBOLIC] for the monotonicity of L_K
% (results/N5_asymptotics.py check (v)); the rest is arithmetic.
%
% R1 says predictive greedy achieves >= L_K(eta) on every instance with error
% <= eta (indeed with eta^path <= eta, a weaker hypothesis).  The theorem says no
% bounded-query algorithm exceeds L_K(hateta).  Since hateta = Phi^{-1}(eta) <= eta
% and L_K is strictly DECREASING in its argument (check (v)), we get
%       L_K(hateta) >= L_K(eta),
% so the two statements are compatible and the theorem is strictly weaker than a
% matching bound.  The residual gap at fixed K is
%       L_K(hateta) - L_K(eta) ~ delta * e^{-1/eta}/eta ~ 2(c+1)(1-1/eta) e^{-1/eta}/(eta K)
% (using the first-branch asymptotics delta ~ 2 tau(1-1/eta)/K, check (i)),
% i.e. Theta(c/K); it vanishes iff c = o(K).  The word "tight" is therefore NOT
% used anywhere in this file (GLOSSARY.md: "tight" must name one of three
% meanings; none of them applies here at finite K).
% Sanity numbers (informational block of results/N5_asymptotics.txt), eta = 2,
% tau = 1:
%    K =  8:  L_K(eta) 0.403281   L_K(hateta) 0.448775   V_j 0.421599   U_K 0.424170
%    K = 16:  L_K(eta) 0.398290   L_K(hateta) 0.418997   V_j 0.407614   U_K 0.408230
%    K = 32:  L_K(eta) 0.395859   L_K(hateta) 0.405758   V_j 0.400560   U_K 0.400711
% HONEST READING OF THIS TABLE.  At every tabulated (K, eta, tau) the constant
% L_K(hateta) proved here is LARGER than U_K(eta) and than V_j(eta), i.e. at
% finite K this theorem is weaker than both of the greedy-side reference curves
% and weaker than the LP limit of the same technique (D2).  Its content is
% therefore asymptotic: all four curves converge to 1 - e^{-1/eta} as K -> oo
% (Corollary N5.8), and only in that limit does the hardness side meet R1.
% The tau = 2 rows are worse still (K = 8, eta = 2: L_K(hateta) = 0.512), and
% the K = 4, tau = 2 rows violate hypothesis (H) outright.


% ===========================================================================
% 6.  Vacuity check (CLAUDE.md, mandatory; one entry per qualifier)
% ===========================================================================
% Procedure: delete the qualifier, or replace it by its opposite, and re-read
% the statement.  Unchanged content => delete the qualifier.  Changed => one
% sentence saying what changed, recorded here and reflected in the text.
%
% V1  "deterministic".
%     Opposite (randomized): the conclusion changes shape -- one cannot fix a
%     single bad O, and the bound acquires the additive slack epsilon_n of
%     Corollary N5.7.  KEEP; the flipped version is stated as that corollary.
%     Note that only averaging over the random string is used, so writing
%     "by Yao's principle" is decoration: the minimax direction is never needed.
%     ACTION: say "averaging over the random string", not "by Yao's principle".
%
% V2  "adaptive".
%     Deleting it changes nothing: "algorithm" already allows adaptivity, and
%     the canonical-transcript device of Lemma N5.3 fixes the query sequence at
%     zero cost.  ACTION: the word is DELETED from the theorem statement; the
%     observation that adaptivity is free lives in this comment only.
%
% V3  "at most Q = n^c queries" (poly-query).
%     Opposite (unbounded queries): the statement becomes FALSE, not merely
%     unproven.  On this family G_O is maximized over K-sets uniquely at O
%     (G_O(O) = 1, and every other K-set has G_O < 1), so exhaustive search over
%     all C(n,K) K-subsets recovers O and achieves ratio 1.  KEEP; this is
%     exactly why GLOSSARY.md forbids the bare phrase "any algorithm": range (i)
%     of that entry gives 1/eta (R2), range (ii) gives this theorem.
%
% V4  "each query on a set of size at most K".
%     Opposite (arbitrary query size): the argument breaks and the family is
%     known to fail.  For |S| = m the hypergeometric mean of y is mK/n, so
%     "y <= tau" stops being the high-probability event once m >> n/K; the
%     correct event is then the concentration band |y - K|S|/n| <= tau, and T2
%     Conclusion 2 exhibits a structural certificate showing THIS F is infeasible
%     under that band for every delta and every n.  KEEP.
%     Quantitatively the restriction can be relaxed to |S| <= M provided
%     Q C(M,tau+1)(K/n)^{tau+1} + MK/n < 1/2, which holds for M up to about
%     n^{2/(c+2)}/K; it cannot be relaxed to M = Theta(n).
%
% V5  "outputs a set T with |T| <= K".
%     Deleting it changes the problem (the cardinality constraint IS the
%     problem) and also breaks the quantitative step E|T_0 cap O| = |T_0|K/n
%     <= K^2/n.  KEEP.
%
% V6  "monotone submodular".
%     Truth value unchanged if replaced by "arbitrary set function" (a larger
%     instance class only makes hardness easier), but the CONTENT changes: the
%     point of the qualifier is that the hardness already holds inside the class
%     where R1 guarantees L_K(eta), so the two constants are comparable.
%     KEEP, with that sentence in the text.
%
% V7  "single-element error at most eta".
%     Opposite (all-pairs error at most eta): NOT SUPPORTED.  The LP bands of
%     T2 constrain adjacent pairs only, so the family's all-pairs error was
%     never verified; since {all-pairs <= eta} is a strictly smaller class,
%     hardness under all-pairs error would be a strictly stronger claim.
%     KEEP the single-element wording and flag the all-pairs version as open
%     (honest declaration D5).  Related: eta^path <= single-element eta, so R1
%     applies to these instances, which is what makes the consistency check of
%     Section 5 meaningful.
%
% V8  "n >= 4K^{c+2}".
%     Deleting it makes the union bound of Lemma N5.2 false and the theorem
%     unproven; small n is not a technicality here (N4 Section 3 shows the
%     related LP value genuinely drops when the ground set is too small to hold
%     the construction's chain).  KEEP; the threshold is sufficient, not
%     necessary, and the CONSTANT L_K(hateta) does not depend on n.
%
% V9  "K >= 2 tau" (= 2c+2; it implies the weaker K >= c+2).
%     Deleting it makes delta undefined at tau = K (K - tau = 0) and breaks the
%     branch-(vi.b) monotonicity that Lemma N5.5 needs (theta K > 2 tau - 1).
%     KEEP: for a fixed query exponent c the theorem says nothing about
%     K < 2c+2, and it degrades smoothly as tau approaches K/2 (delta =
%     Theta(tau/K) = Theta(1) there).
%
% V10 "there is a set O" (existential / worst case).
%     Opposite ("for every O"): FALSE, since an algorithm may hard-code one O.
%     KEEP.
%
% V11 "tau = c+1".
%     Replacing it by tau = c leaves the union bound valid (margin 1/n instead
%     of 1/n^2) and DECREASES delta, i.e. strengthens the conclusion.  So this
%     qualifier is a choice, not a necessity; it is kept only because TASKS2.md
%     fixes it and because tau = 0 is degenerate.  RECORDED as a known
%     non-optimality (honest declaration D4).
%
% V12 "information-theoretic" (used in the surrounding prose, not in the
%     statement).  Per GLOSSARY.md it means: the bound depends only on the
%     number of queries and uses no complexity assumption.  Deleting it would
%     leave readers free to read a poly-TIME claim into the theorem, which is
%     false -- nothing here is a computational lower bound.  KEEP, with the
%     GLOSSARY definition given at first use.
%
% V13 "at error eta" versus the eta-hat inside the bound.
%     Not a deletable qualifier but the most error-prone phrase in the
%     statement: the promise is "error at most eta", while the hardness constant
%     is evaluated at the strictly smaller hateta = Phi^{-1}(eta).  The direction
%     is forced (a hardness instance must lie INSIDE the promised class), and it
%     is what makes the bound weaker than L_K(eta).  Flipping it -- reading the
%     statement as "the family attains error eta and forces L_K(eta)" -- is the
%     single most likely misreading, and it is FALSE.  See honest declaration D3.
%
% V14 "any algorithm" (contribution-sentence phrasing).
%     Bare use is banned by GLOSSARY.md.  ACTION: always write "any
%     deterministic algorithm that makes n^c queries on sets of size at most K".
%
% V15 "tight".
%     Not used, and must not be added: at finite K the theorem leaves a
%     Theta(c/K) gap (Section 5), and the K -> infinity matching statement is
%     conditional on a [CONJECTURE].
%
% V16 hypothesis (H), "Phi(1) <= eta".
%     Deleting it makes the theorem FALSE as stated, not merely unproven: the
%     design parameter would fall below 1, and for theta < 1 Lemma N5.4 itself
%     fails (its two "<= 1" edge bounds use theta >= 1), quite apart from the
%     fact that an error level below 1 is not a legal value of eta = eta_u eta_o.
%     KEEP.  This qualifier is MISSING from the target statement in TASKS2.md
%     and was added here; results/N5_delta_at_etahat.txt marks the two points
%     where it fails as SKIP(H).
%
% V17 "eta > 1".
%     Opposite (eta = 1, perfect predictions): a = (K-1)/K, so the first branch
%     ((K-1)/K)^tau K/(K-tau) equals 1 exactly at tau = 1 and is > 1 for every
%     tau >= 2 (its log is (tau^2 - tau)/(2K^2) + O(1/K^3) > 0), while the second
%     branch a^{1-tau} = (K/(K-1))^{tau-1} > 1 for tau >= 2.  So delta(1) = 0 and
%     (H) holds with equality only when tau = 1; for tau >= 2 (H) FAILS at every
%     K.  Since (H)'s sufficient form K >= tau(1 + 2/ln eta) blows up as eta -> 1,
%     the theorem is vacuous in a neighbourhood of eta = 1 whenever c >= 1.
%     KEEP eta > 1 and state that the consistency regime (LAA sense,
%     GLOSSARY.md) is outside its reach.


% ===========================================================================
% 7.  Honest declarations
% ===========================================================================
% D1  Verification scope of the delta formula, and the correction to T2.
%     (a) T2's own scope [VERIFIED-LP, finite instances]: full-lattice LP at
%         (n,K) in {(8,3),(10,3),(10,4),(12,4)}; symmetrized grid at
%         K in {3,4,6,8,12,16,24,32}, n = 4K; eta in {3/2, 2, 3}; tau in {1,2};
%         error splits (sqrt eta, sqrt eta) and (eta, 1).
%     (b) T2's REPORTED CLOSED FORM IS INCOMPLETE.  The correct expression is
%         the max of Lemma N5.4; the T2 form is its first branch.  The two agree
%         on every T2 test point (that is why the gap was invisible), but the
%         theorem evaluates the family at a design parameter strictly below eta,
%         and there they differ.  results/N5_delta_at_etahat.py solves 40 such
%         points: LP = max-form at 40/40, LP != T2-form at 8/40 (listed in the
%         Lemma N5.4 comment; worst case a factor 1.93 in 1 + delta).
%         ACTION: amend results/T2_summary.md Conclusion 1 and
%         RESEARCH_STATE.md R11(a).  Not done here (N5 writes results/N5_* only).
%     (c) Remaining [CONJECTURE]: the max-form for general (K, tau, theta).
%         Both scripts together cover tau in {1,2} only, i.e. c in {0,1}
%         (Q = O(1) and Q = n).  For c >= 2 the theorem's constant rests on a
%         [CONJECTURE], and the second branch a^{1-tau} grows with tau, so the
%         untested regime is exactly the one where the correction matters most.
%     (d) The extrapolation gap flagged in an earlier draft of this file (T2
%         verified the formula only at theta in {3/2,2,3}, while the theorem
%         needs it at theta = hateta) is CLOSED at 40 points by
%         results/N5_delta_at_etahat.py, which runs the T2 LP oracle at exactly
%         those hateta values.
%
% D2  The F of this theorem is FIXED and is not the best available.
%     The theorem uses the R9 candidate F(x,y) = hateta^{-1/2}(1 - hata^x(1-y/K)),
%     which is the R7/U_K shape.  N4 (results/N4_hardness_construction.md)
%     shows that when F is allowed to vary, the same LP -- with delta = 0 and
%     the same y <= tau balanced region -- has optimum equal to
%     V_j(eta) + (tiny excess), where V_j is R10's closed form rho_K^{LP}, and
%     it exhibits an explicit optimal (F, G) [EXACT rational feasibility at 42
%     (K, eta) points, tau = 1 only].  Numerically V_j < L_K(hateta)
%     (K = 8, eta = 2: 0.4216 vs 0.4488), so feeding N4's family into Lemmas
%     N5.1--N5.3 would give a STRICTLY STRONGER theorem, and with hateta = eta
%     (no inflation at all, since N4's construction is feasible at delta = 0).
%     That is not done here because N4 tested only tau = 1 (i.e. c = 0), its
%     general-K submodularity is 42-point exact-rational rather than symbolic,
%     and its closed form is optimal only for eta <= K-1.  So: the bound of this
%     theorem is NOT the strongest one this technique is known to support.
%
% D3  "error exactly eta" versus "error at most eta(1+delta)" -- getting the
%     direction right.  A hardness instance must lie INSIDE the class the
%     algorithm is promised, so what must be checked is
%     (family error) <= eta, never the reverse.  The family with design
%     parameter theta has single-element error exactly Phi(theta) =
%     theta(1 + delta(theta)) (Lemma N5.4).  Therefore:
%      (a) The theorem takes hateta = Phi^{-1}(eta), so the family's error is
%          EXACTLY eta -- no slack is thrown away, and this is the strongest
%          member of the family admissible at error eta (Phi increasing, L_K
%          decreasing).  Phi^{-1} is a one-dimensional root find, not a formula.
%      (b) TASKS2.md's explicit hateta = eta/(1 + delta(eta)) is a LOWER bound
%          for Phi^{-1}(eta) only when delta(hateta) <= delta(eta).  With the
%          corrected max-form delta is NOT monotone in theta (its first branch
%          increases, its second decreases), so this must be checked rather than
%          assumed; it held at all 40 points of results/N5_delta_at_etahat.py
%          (column E) and the resulting bound L_K(eta/(1+delta(eta))) is weaker
%          than the theorem's.  This is the one place where an earlier draft of
%          this file used a monotonicity argument that the corrected delta does
%          not support.
%      (c) Either way the constant is L_K(hateta) > L_K(eta), so the theorem
%          does NOT match R1 at finite K (Section 5); the mismatch is the price
%          of the concentration argument, not an artefact of the bookkeeping.
%
% D4  Known non-optimalities of the statement itself.
%     (i)  tau = c+1 is not optimal (V11): tau = c also suffices for the union
%          bound and gives a smaller delta, hence a stronger bound.
%     (ii) n >= 4K^{c+2} is a sufficient, not a necessary, threshold (V8).
%     (iii) hypothesis (H) is stated with a crude sufficient form
%          K >= tau(1 + 2/ln eta); the sharp condition is Phi(1) <= eta.
%     (iv) The randomized corollary's epsilon_n is not optimized.
%
% D5  Model and method caveats.
%     (i)  All LP evidence comes from the symmetrized (x,y)-grid LP;  T2
%          cross-checked symmetrization against the full lattice only at
%          (n,K) in {(8,3),(10,3),(10,4),(12,4)}.  The family here is explicit,
%          so symmetrization is needed only for the MINIMALITY of delta, not for
%          feasibility.
%     (ii) All-pairs error unverified (V7).
%     (iii) Monotonicity/submodularity of F: the (x,y) difference structure is
%          the R7 family up to the parameter name, hence [VERIFIED-SYMBOLIC] for
%          general K via results/T5_symbolic.py, and [VERIFIED-LP] on the full
%          lattice for n <= 12 via T2.  There is no separate general-n symbolic
%          check in this file.
%     (iv) OPT = f(O) and the ratio identity (F3), (F4) are [VERIFIED-LP] by
%          enumeration at the T2 test points and, at the theorem's own hateta
%          values, by results/N5_delta_at_etahat.py (columns C and D).
%     (v)  Queries are to ftilde only; f is inaccessible (paper's model).  The
%          theorem counts ftilde-queries and nothing else.
%     (vi) The randomized corollary bounds min over the FAMILY of Section 0, not
%          over all monotone submodular instances with error <= eta; that is the
%          usual (and sufficient) form of an averaging lower bound, but it should
%          be phrased that way in the paper rather than as "no randomized
%          algorithm can do better than L_K(hateta)".
%
% ===========================================================================
% Dependency summary (script / file per step)
% ===========================================================================
%   Section 0 family .......... RESEARCH_STATE.md R9; results/T2_hardness_grid.py
%   Lemma N5.1 (conc.) ........ none (hand)                [HAND-PROOF-UNREVIEWED]
%   Lemma N5.2 (union) ........ none (hand)                [HAND-PROOF-UNREVIEWED]
%   Lemma N5.3 (valuation) .... facts (F1)-(F4) from R11(a) / T2_summary.md
%                               ([VERIFIED-LP finite]); argument [HAND-PROOF-UNREVIEWED]
%   Lemma N5.4 (delta) ........ hand derivation [HAND-PROOF-UNREVIEWED] +
%                               results/N5_delta_at_etahat.py (40 PASS / 0 FAIL,
%                               8 points refuting the T2 form) [VERIFIED-LP finite];
%                               supersedes results/T2_summary.md Conclusion 1 and
%                               results/T2_hardness_grid.py:261
%   Lemma N5.5 (Phi invertible) results/N5_asymptotics.py (iv),(vi) [VERIFIED-SYMBOLIC]
%   Theorem N5.6 .............. assembly                    [HAND-PROOF-UNREVIEWED]
%   Corollary N5.7 (random) ... assembly                    [HAND-PROOF-UNREVIEWED]
%   Corollary N5.8 (asympt.) .. results/N5_asymptotics.py (i),(ii) [VERIFIED-SYMBOLIC]
%   Section 5 consistency ..... results/N5_asymptotics.py (v)  [VERIFIED-SYMBOLIC]
%   D2 comparison numbers ..... results/N5_asymptotics.txt informational block;
%                               results/N4_hardness_construction.md Sections 3, 6
%
% Reproduce:
%   python3 results/N5_asymptotics.py      -> results/N5_asymptotics.txt   (15/15, exit 0)
%   python3 results/N5_delta_at_etahat.py  -> results/N5_delta_at_etahat.txt
%                                             (40 PASS / 0 FAIL / 2 SKIP, exit 0)
